\chapter{Binary Detection Results}
\label{ch:binary}

\section{Headline Numbers}

Reframing the task as ``benign versus any-attack'' recovers most of the
performance lost in the multi-class collapse. Random forest reaches
F1 = 0.86 on the day split, within 0.14 of its stratified score of 0.996.
XGBoost and LightGBM trail random forest on the day split (F1
0.74--0.76) despite higher stratified scores, because their decision
boundaries over-fit the Mon--Wed attacks.

\begin{table}[H]
    \centering
    \caption{CICIDS2017 binary F1 at the Youden-J threshold.}
    \label{tab:binary-headline}
    \small
    \begin{tabular}{lcccc}
        \toprule
        \textbf{Model} & \textbf{Stratified F1} & \textbf{Day F1} & \textbf{Day Threshold} & \textbf{Day ROC-AUC} \\
        \midrule
        LogReg & 0.896 & 0.681 & 0.0627 & 0.981 \\
        Random Forest & 0.996 & \textbf{0.859} & 0.0090 & 0.940 \\
        XGBoost & 0.997 & 0.757 & 0.000323 & 0.972 \\
        LightGBM & 0.998 & 0.744 & 0.000178 & 0.957 \\
        \bottomrule
    \end{tabular}
\end{table}

\interpretation{The thresholds in column four are striking: tree
ensembles need to operate at probabilities below $10^{-3}$ on the day
split. Reporting F1 at the canonical 0.5 threshold would yield F1
$\approx 0.007$ for XGBoost. Threshold transfer is not optional --- it
is a first-class deployment concern.}

\section{Operating Points and Calibration}
\label{sec:operating-points}

For deployment, the question is not which threshold maximises F1 but
which threshold meets a fixed false-positive budget. A SOC can absorb a
finite number of false alerts per analyst per shift, so the
operationally relevant metric is recall at a fixed FPR.
Table~\ref{tab:operating-points} reports recall at three FPR budgets
(0.01\%, 0.1\%, 1\%) per model on the binary day split, alongside the
Expected Calibration Error (ECE) that explains the pattern.

\begin{table}[H]
    \centering
    \caption{Recall at fixed FPR budgets on the binary day split
    ($n = 516{,}643$, attack rate $0.2541$). ECE is the Expected
    Calibration Error with 15 equal-width bins; lower is better.}
    \label{tab:operating-points}
    \small
    \begin{tabular}{lcccc}
        \toprule
        \textbf{Model} & \textbf{ECE $\downarrow$} & \textbf{Recall @ 0.01\% FPR} & \textbf{Recall @ 0.1\% FPR} & \textbf{Recall @ 1\% FPR} \\
        \midrule
        LogReg & 0.073 & 0.393 & 0.394 & 0.605 \\
        Random Forest & 0.178 & 0.000 & \textbf{0.621} & 0.634 \\
        XGBoost & 0.252 & 0.001 & 0.031 & \textbf{0.637} \\
        LightGBM & 0.249 & 0.001 & 0.461 & 0.620 \\
        \bottomrule
    \end{tabular}
\end{table}

A 0.1\% FPR budget on a two-million-flow day produces around 2{,}000
false alerts, which is the regime where a small SOC can absorb the alert
volume without burning analyst hours; the 1\% budget is closer to a
large enterprise with more analysts on the queue. Random forest
dominates the strict end of the curve and XGBoost catches up only as the
FPR budget loosens.

\begin{figure}[H]
    \centering
    \includegraphics[width=0.8\textwidth]{operating_points.png}
    \caption{Recall at fixed FPR budgets per model on the binary day
    split. At a 1\% FPR budget all four models recover similar recall
    (60--64\%); at the tighter 0.1\% FPR budget, random forest still
    catches 62\% of attacks while XGBoost drops to 3\%.}
    \label{fig:op-points}
\end{figure}

Calibration explains why the operating-point picture varies so much
between models. Logistic regression on the day split is nearly perfectly
calibrated (ECE = 0.073); tree ensembles are markedly less so (RF 0.18,
XGB and LGB 0.25). Poor calibration is exactly why the canonical 0.5
threshold fails for tree ensembles under temporal shift, and why
Youden's J on validation is needed to recover sensible operating points.

\section{Statistical Significance}
\label{sec:mcnemar}

Two statistical tests support the model ranking. McNemar's test with
Yates continuity correction on every pair of models on the binary
day-split test set rejects the null of equal error rates at $p < 10^{-4}$
for all six pairs, including the closest (XGBoost versus LightGBM,
$\chi^2 = 520.6$). Bootstrap 95\% confidence intervals (with
$B = 1000$ resamples on $n = 516{,}643$ test flows) are correspondingly
narrow, typically about $\pm 0.001$ on F1, and confirm the same ordering.
Random forest has the highest F1 with non-overlapping precision/recall
confidence intervals against all other models.
